% !TEX program = xelatex
% AIPL Runtime Feature Matrix — 7 implementations compared

\documentclass[10pt,a4paper,landscape]{article}

\usepackage[a4paper,landscape,margin=15mm]{geometry}
\usepackage{xeCJK}
\usepackage{fontspec}
\usepackage{amsmath,amssymb}
\usepackage{xcolor}
\usepackage{booktabs}
\usepackage{tabularx}
\usepackage{longtable}
\usepackage{array}
\usepackage{makecell}
\usepackage{hyperref}
\usepackage{enumitem}

\setCJKmainfont{Hiragino Mincho ProN}
\setCJKsansfont{Hiragino Sans}
\setCJKmonofont{Hiragino Sans}
\setmonofont{Menlo}[Scale=0.85]

\hypersetup{
  colorlinks=true,
  linkcolor=black,
  urlcolor=blue!60!black,
  pdftitle={AIPL Runtime Feature Matrix},
  pdfauthor={Yasushi Kodama}
}

% Status cells — colored backgrounds for readability
\newcommand{\Yes}{\textcolor{green!50!black}{\textbf{\checkmark}}}
\newcommand{\No}{\textcolor{red!70!black}{$\times$}}
\newcommand{\Part}{\textcolor{orange!80!black}{$\triangle$}}
\newcommand{\NA}{\textcolor{gray}{--}}

% Column types for tight matrix
\newcolumntype{C}{>{\centering\arraybackslash}p{0.9cm}}
\newcolumntype{F}{>{\raggedright\arraybackslash}p{6.5cm}}
\newcolumntype{N}{>{\centering\arraybackslash}p{0.6cm}}

\title{\textbf{AIPL Runtime Feature Matrix} \\
  {\large 7 implementations compared}}
\author{Yasushi Kodama}
\date{2026-05-13}

\begin{document}
\maketitle

\section*{Overview}

AIPL (Actor-based Intelligent Parallel Language, formerly ABCL/c+)
has grown into a \textbf{seven-runtime project}.  This document
provides a side-by-side comparison of every language and
infrastructure feature across all implementations, grounded in
source surveys.

\paragraph{Legend.}
\Yes\ = full support, \Part\ = partial, \No\ = not supported,
\NA\ = not applicable.

\paragraph{The seven runtimes.}

\begin{tabularx}{\linewidth}{l X}
  \toprule
  \textbf{Short name} & \textbf{Source} \\
  \midrule
  Python (annotated)        & \texttt{src/python-aipl/} \\
  Python (inferred)         & \texttt{src/python-aipl-inferred/} \\
  OCaml                     & \texttt{src/*.ml}, \texttt{\_build/.../repl\_thread.exe} \\
  JS-OCaml (server)         & \texttt{src/app.js}, \texttt{console\_server.js}, \texttt{ide.js} + OCaml \texttt{web\_gateway} \\
  JS-Browser (serverless)   & \texttt{src/browser-abcl/} \\
  JS-Node (server)          & \texttt{src/node-aipl-server/} (Node HTTP server wrapping browser-abcl) \\
  C (abcl2c)                & \texttt{src/abcl2c.ml} codegen + \texttt{src/abcl\_gui\_runtime.c} \\
  \bottomrule
\end{tabularx}

\paragraph{Column abbreviations in matrices below.}
\textbf{Py-A}~=~Python annotated, \textbf{Py-I}~=~Python inferred,
\textbf{OCaml}, \textbf{JS-O}~=~JS via OCaml backend,
\textbf{JS-B}~=~JS in browser, \textbf{JS-N}~=~JS in Node server,
\textbf{C}~=~abcl2c codegen.

\paragraph{Inheritance relationships.}

\begin{itemize}[leftmargin=*,itemsep=0pt]
  \item \textbf{JS-OCaml} is a thin HTTP client over the OCaml
    backend, so its feature set inherits OCaml's (plus the new
    \texttt{/api/typecheck} endpoint).
  \item \textbf{JS-Node} wraps the same browser-abcl parser /
    type-checker / interpreter inside a Node.js HTTP server.
    Runtime feature set therefore mirrors JS-Browser.
  \item \textbf{Python (inferred)} imports parser, interpreter, AI,
    remote, and dashboard from \texttt{../python-aipl/}.  All
    \emph{runtime} features match Python (annotated); only the
    \emph{static type-checker} layer differs (full
    Hindley-Milner instead of Phase 11+ annotation-driven).
\end{itemize}

\clearpage

\section*{Sample programs}

Number of \texttt{.abcl} sample programs reachable by each runtime
(core directory + AI integration + remote-actor demos):

\begin{tabularx}{\linewidth}{l C C C C C C C}
  \toprule
  & \textbf{Py-A} & \textbf{Py-I} & \textbf{OCaml} & \textbf{JS-O} & \textbf{JS-B} & \textbf{JS-N} & \textbf{C} \\
  \midrule
  core samples           & 33   & 33   & 57   & 57   & 5    & 5    & 57   \\
  AI integration samples & 11   & 11   & 8    & 8    & 0    & 0    & 8    \\
  remote-actor samples   & 8    & 8    & 1    & 1    & 0    & 0    & 1    \\
  \midrule
  \textbf{total}         & \textbf{52} & \textbf{52} & \textbf{66} & \textbf{66} & \textbf{5} & \textbf{5} & \textbf{66} \\
  \bottomrule
\end{tabularx}

\smallskip

\noindent
Cross-cutting / not tied to a specific runtime:
\begin{itemize}[leftmargin=*,itemsep=0pt]
  \item \texttt{aipl-self-host/}: 48 AIPL-in-AIPL self-host programs
  \item \texttt{docker/cross/samples/}: 4 cross-language interop samples
\end{itemize}

Note that the C runtime processes the same \texttt{.abcl} files as
OCaml via \texttt{abcl2c}; of the 66 reachable samples, 48 currently
pass the \texttt{abcl2c} smoke test (the remaining 9 in \texttt{abclc/}
have pre-existing type ambiguities surfaced by our hard-fail policy).

\clearpage

\section*{Core language}

\begin{longtable}{F C C C C C C C}
  \toprule
  \textbf{Feature} & \textbf{Py-A} & \textbf{Py-I} & \textbf{OCaml} & \textbf{JS-O} & \textbf{JS-B} & \textbf{JS-N} & \textbf{C} \\
  \midrule
  \endhead

  method injection (\texttt{add\_method}/\texttt{remove\_method}) & \Yes & \Yes & \No & \No & \No & \No & \No \\
  \texttt{now} synchronous send                                   & \Yes & \Yes & \Yes & \Yes & \Yes & \Yes & \Yes \\
  \texttt{future} async send                                      & \Yes & \Yes & \Yes & \Yes & \Yes & \Yes & \Yes \\
  \texttt{await} future block                                     & \Yes & \Yes & \Yes & \Yes & \Yes & \Yes & \Yes \\
  \texttt{send} fire-and-forget                                   & \Yes & \Yes & \Yes & \Yes & \Yes & \Yes & \Yes \\
  \texttt{become} actor class swap                                & \Yes & \Yes & \Yes & \Yes & \No  & \No  & \No  \\
  \texttt{select} selective receive                               & \Yes & \Yes & \Yes & \Yes & \Yes & \Yes & \No  \\
  top-level functions                                             & \Yes & \Yes & \Part & \Part & \No  & \No  & \Yes \\
  signatures / overloads                                          & \No  & \No  & \Yes & \Yes & \No  & \No  & \No  \\
  dynamic compile (\texttt{compile()})                            & \Yes & \Yes & \No & \No & \No & \No & \No \\
  worker pool / \texttt{DynamicWorkerPool}                        & \Yes & \Yes & \No & \No & \Part & \Part & \No \\
  \bottomrule
\end{longtable}

\section*{Type system}

\begin{longtable}{F C C C C C C C}
  \toprule
  \textbf{Feature} & \textbf{Py-A} & \textbf{Py-I} & \textbf{OCaml} & \textbf{JS-O} & \textbf{JS-B} & \textbf{JS-N} & \textbf{C} \\
  \midrule
  \endhead

  type inference                  & \Part trace & \Yes HM & \Yes HM & \Yes HM & \Yes flow & \Yes flow & \Yes HM+spec \\
  type annotations                & \Yes        & \Part ignored & \No & \No & \No & \No & \No \\
  records \texttt{\{a:int, b:str\}} & \Yes  & \Yes & \Part type only & \No & \No & \No & \Part type only \\
  tuples \texttt{(1, "a")}          & \Yes  & \Yes & \No  & \No  & \No  & \No  & \No  \\
  arrays (typed, multi-dim)         & \Yes  & \Yes & \Yes & \Yes & \Part & \Part & \Part \\
  generics on functions             & \No   & \Yes Forall & \Part Forall & \Part & \No & \No & \No \\
  \bottomrule
\end{longtable}

\section*{Phase 11+ advanced features}

Python (inferred) runs HM only, so the annotation-driven Phase 11+
static checks that Python (annotated) performs are not re-implemented.
Programs still execute at runtime via the shared interpreter, but
the static guarantees from Phase 12 / 14 / 15 / 16 are lost.
Phase 13 (channels) and Phase 17 (\texttt{scope}) are runtime features,
so they remain available.

\begin{longtable}{F C C C C C C C}
  \toprule
  \textbf{Feature} & \textbf{Py-A} & \textbf{Py-I} & \textbf{OCaml} & \textbf{JS-O} & \textbf{JS-B} & \textbf{JS-N} & \textbf{C} \\
  \midrule
  \endhead

  channels (CSP-style) — Phase 13          & \Yes & \Yes & \No & \No & \No & \No & \No \\
  linear types / use-after-move — Phase 14 & \Yes & \No  & \No & \No & \No & \No & \No \\
  owned / pub fields — Phase 15            & \Yes & \No  & \No & \No & \No & \No & \No \\
  effects \texttt{!\{fs,ai,net,mut\}} — Phase 12 & \Yes & \No  & \No & \No & \No & \No & \No \\
  structured concurrency \texttt{scope} — Phase 17 & \Yes & \Yes & \No & \No & \No & \No & \No \\
  transient cast at any-boundary — Phase 16 & \Yes & \No  & \No & \No & \No & \No & \No \\
  \bottomrule
\end{longtable}

\clearpage

\section*{Networking, AI, infrastructure}

\begin{longtable}{F C C C C C C C}
  \toprule
  \textbf{Feature} & \textbf{Py-A} & \textbf{Py-I} & \textbf{OCaml} & \textbf{JS-O} & \textbf{JS-B} & \textbf{JS-N} & \textbf{C} \\
  \midrule
  \endhead

  remote actors (HTTP)                          & \Yes & \Yes & \Part no WS & \Part & \No & \Part srv only & \Yes via OCaml \\
  WebSocket                                     & \No & \No & \No & \No & \No & \No & \No \\
  AI integration (\texttt{ai\_call})            & \Yes stream/img & \Yes & \Yes & \Yes & \Yes mock & \Yes mock & \Yes \\
  AI governance (budget/concurrent/fallback)    & \Yes & \Yes & \Yes & \Yes & \No & \No & \Yes \\
  HMAC-signed remote send                       & \Yes & \Yes & \No & \No & \No & \No & \Yes \\
  persistent actor state                        & \Yes & \Yes & \No & \No & \No & \No & \No \\
  live dashboard (SSE)                          & \Yes & \Yes & \Part polling & \Part & \No & \No & \Yes \\
  \texttt{/api/typecheck} JSON endpoint         & \No & \No & \Yes & \Yes & \No & \Yes & \NA \\
  \texttt{/api/run} JSON endpoint               & \No & \No & \No & \No & \No & \Yes & \NA \\
  \bottomrule
\end{longtable}

\section*{C-version-specific codegen targets}

\begin{tabularx}{\linewidth}{X C}
  \toprule
  \textbf{Target} & \textbf{C} \\
  \midrule
  C + pthread standalone binary              & \Yes \\
  C + SDL2 GUI binary (1178-line runtime)    & \Yes \\
  Xinu embedded-OS target (\texttt{--xinu})  & \Yes \\
  Python target (\texttt{--python})          & \Yes \\
  \bottomrule
\end{tabularx}

\clearpage

\section*{Key observations}

\subsection*{Python (annotated) is the most feature-complete runtime}

Of the 30 surveyed features, Python (annotated) ticks roughly 28.
It owns the full Phase 11--17 stack (channels, linear types, owned
fields, effects, structured concurrency, transient casts), method
injection, dynamic compile, persistent actor state, and full AI
integration (including streaming and images).

\subsection*{Python (inferred) is Python's HM-typed twin}

Runtime features match Python (annotated) exactly because the
interpreter is shared.  The trade is on the static side: the Phase
11+ annotation-driven checks are replaced with a Hindley-Milner
inference pass.  Inferred types cross-verified against OCaml: on
shared samples (\texttt{Hello.abcl}, \texttt{counter.abcl}) the two
runtimes produce identical type schemes.

\subsection*{OCaml is the canonical core}

Implements roughly 15 of the surveyed features.  Phase 11+ exist
only as \texttt{.abcl} design-document samples, not implemented.
Solid \texttt{become} / \texttt{select} / HM inference; remote is
HTTP only (no HMAC, no WebSocket).

\subsection*{JS-OCaml $\approx$ OCaml}

Thin HTTP client over the OCaml backend, so features inherit.  The
only differentiator is the \texttt{/api/typecheck} endpoint added
in commit \texttt{5227f87}.

\subsection*{JS-Browser is the minimal implementation}

Around 8 of the 30 features.  Basic actor model plus our
newly-added flow-sensitive type inference.  No \texttt{become}, no
channels, no remote.  AI integration exists but is mock-only.

\subsection*{JS-Node (Node-server, newest)}

Wraps the browser-abcl parser, type-checker, and interpreter
inside a Node.js HTTP server.  Exposes \texttt{/api/typecheck} (the
same JSON shape as OCaml's endpoint) and \texttt{/api/run}
(executes and returns stdout) — the only runtime with the latter.
Runtime feature set mirrors browser-abcl; the value-add is having
both the type-checker and the interpreter reachable over HTTP
without bringing up either OCaml or a browser.

\subsection*{C version --- surprisingly strong on infrastructure}

About 17 of 30.  \texttt{become} and \texttt{select} are not
implemented (codegen emits \texttt{/* unsupported */}), but it
inherits OCaml's HMAC, SSE, AI governance, and adds \textbf{three
additional codegen targets} (SDL2, Xinu, Python).  Fields,
parameters, and locals are specialized to native C types.

\section*{Method injection --- a Python-family exclusive}

Both Python variants use mutable method dispatch tables that can
be mutated at runtime via \texttt{add\_method} /
\texttt{remove\_method}.  Other runtimes substitute \texttt{become}
(whole-class swap) as the closest equivalent.  \texttt{JS-Browser}
and \texttt{JS-Node} lack even \texttt{become} (pure actor
execution).

\section*{Type inference cross-verification}

For shared samples (\texttt{abclc/Hello.abcl},
\texttt{abclc/counter.abcl}), the three HM-based runtimes
(\textbf{OCaml}, \textbf{Python-inferred}, \textbf{C}) produce
identical inferred types:

\begin{verbatim}
Hello:   count : float
         init  : (int) -> unit
         greet : () -> unit
         inc   : () -> unit

Counter: count : float
         inc   : () -> unit
         dec   : (float) -> unit    <- monomorphized via call site
\end{verbatim}

The flow-sensitive variants (\textbf{JS-Browser},
\textbf{JS-Node}) reach the same answers on shared samples within
the limits of their algorithm — field types match, but method
signature schemes are not produced the same way.

\section*{Phase 11+ implementation gap}

The full Phase 11--17 stack is implemented only in Python
(annotated).  This represents AIPL's language-evolution frontier;
Python (annotated) serves as the research runtime.  Python
(inferred) keeps the runtime side but trades the static checks for
HM-style polymorphic inference --- a different research axis
(declarative typing vs. annotated capabilities).

\bigskip
\hrule
\smallskip

\noindent\small
Generated 2026-05-13.  Includes \texttt{python-aipl-inferred}
(commit \texttt{270f291}) and \texttt{node-aipl-server} (commit
\texttt{adeb0b6}).  For source pointers, run \texttt{grep} against
the files listed in each runtime's source column.

\end{document}
